\section{The three approaches compared}
\label{sec:compare}

\subsection{Where the effort went}

\begin{center}
\begin{tikzpicture}[font=\footnotesize\sffamily,
  lay/.style={draw,rounded corners=3pt,minimum width=34mm,minimum height=8mm,align=center,inner sep=3pt},
  a/.style={lay,draw=cGrey,fill=cGreyBg}, h/.style={lay,draw=cRed,fill=cAmberBg,line width=1pt}]
\node[font=\small\bfseries\sffamily] at (0,3.3) {\udp{} 2018};
\node[font=\small\bfseries\sffamily] at (4.5,3.3) {\QED{} 2024};
\node[font=\small\bfseries\sffamily] at (9,3.3) {\sqlsolver{} 2024};
\foreach \x/\t in {0/{}, 4.5/{}, 9/{}} {}
% UDP column
\node[a] at (0,2.5) {SQL $\to$ U-expression};
\node[h] at (0,1.5) {\textbf{U-semiring axioms}\\\textbf{keys as identities}};
\node[a] at (0,0.4) {SPNF};
\node[h] at (0,-0.6) {\textbf{syntactic}\\\textbf{isomorphism}};
% QED column
\node[h] at (4.5,2.5) {\textbf{Q-expr over $\Nbar$}\\\textbf{NULLs, QOp/HOp}};
\node[h] at (4.5,1.5) {\textbf{PK/FK as rewrites}\\\textbf{+ global formulas}};
\node[h] at (4.5,0.4) {\textbf{stabilise (SMT}\\\textbf{congruence)}};
\node[h] at (4.5,-0.6) {\textbf{unify + ambient}\\\textbf{assumption, SMT}};
% SQLSolver column
\node[h] at (9,2.5) {\textbf{extended U-expr}\\\textbf{concrete aggregates}};
\node[a] at (9,1.5) {constraints\\(inherited)};
\node[a] at (9,0.4) {vectorise sums};
\node[h] at (9,-0.6) {\textbf{\liastar{} $\to$ \lia{}}\\\textbf{$\to$ SMT}};
\node[align=left,font=\scriptsize\sffamily,text=cGrey] at (4.5,-1.7)
  {\textcolor{cRed}{\rule{3mm}{2mm}} highlighted = where that system innovated \qquad
   \textcolor{cGrey}{\rule{3mm}{2mm}} grey = inherited or routine};
\end{tikzpicture}
\end{center}

\noindent The picture says it: \QED{} rebuilt the \emph{modelling} and the \emph{search};
\sqlsolver{} rebuilt the \emph{translation of aggregates} and the \emph{arithmetic backend}.
They are close to orthogonal.

\subsection{Capability matrix}

\begin{table}[h]\centering\small
\begin{tabular}{@{}lccc@{}}
\toprule
 & \udp{} & \QED{} & \sqlsolver{} \\
\midrule
semantic domain & any U-semiring & $\Nbar$ only & any U-semiring\\
predicate reasoning & congruence closure & SMT & SMT\\
NULL / 3-valued logic & \textcolor{cRed}{no} & \textcolor{cGreen}{full} & partial (from \wetune{})\\
outer joins & \textcolor{cRed}{no} & \textcolor{cGreen}{yes} & partial\\
primary / foreign keys & \textcolor{cGreen}{identities} & \textcolor{cGreen}{rewrite + global formula} & inherited\\
FK cycles terminate & \textcolor{cRed}{no} & \textcolor{cGreen}{yes (lazy)} & --\\
aggregates & uninterpreted & uninterpreted ($\hop$) & \textcolor{cGreen}{concrete}\\
unknown operators & \textcolor{cRed}{unsupported} & \textcolor{cGreen}{$\qop$/$\hop$ first class} & \textcolor{cRed}{unsupported}\\
differing \#summations & \textcolor{cRed}{no} & only inside $\mathcal F_T$ & \textcolor{cGreen}{yes (\liastar{})}\\
\lstinline|ORDER BY| & \textcolor{cRed}{no} & \textcolor{cRed}{no} & \textcolor{cGreen}{yes (separate layer)}\\
bag \lstinline|INTERSECT ALL| & \textcolor{cRed}{no} & \textcolor{cRed}{no (needs monus)} & \textcolor{cGreen}{yes}\\
completeness fragment & UCQ (set \emph{or} bag) & $\mathcal F_T$ (UCQ + any theory $T$) & CQ/UCQ + a group-by fragment\\
refutation & \textcolor{cRed}{no} & \textcolor{cRed}{no} & \textcolor{cRed}{no}\\
\bottomrule
\end{tabular}
\end{table}

\subsection{A gallery of separating examples}

The fastest way to understand three provers is to see what each one alone can do.

\paragraph{(1) Syntactically different but equal predicates.}
\lstinline|WHERE a+b = 0| versus \lstinline|WHERE a = -b|.
\udp{} \textcolor{cRed}{fails} --- congruence closure treats $+$ as uninterpreted.
\QED{} and \sqlsolver{} \textcolor{cGreen}{succeed} --- both call an SMT solver on
predicates. This single limitation is why \emph{everyone} after 2018 uses SMT.

\paragraph{(2) Aggregate pushdown through a key join} (\S\ref{sec:running}).
Needs: primary key as an algebraic fact \emph{and} context-sensitive comparison of the two
aggregate arguments. \QED{} \textcolor{cGreen}{succeeds} via
$S\rightsquigarrow\lam{a}{\sq{S(a)}}$ plus the $\phi_1\vee\phi_2$ ambient assumption.
\udp{} \textcolor{cRed}{fails} on the context-free recursive check (\S\ref{sec:udp-fails}).
\spes{} \textcolor{cRed}{fails} --- no applicable key rule.

\paragraph{(3) Differing numbers of summations} (Example~\ref{ex:splitunion}).
\sqlsolver{} \textcolor{cGreen}{succeeds} via \liastar{}.
\udp{} \textcolor{cRed}{fails} at Alg.~2 line 3 (term counts differ).
\QED{} \textcolor{cGreen}{also succeeds} here --- but only because this query lies in
$\mathcal F_T$, so the \emph{complete-fragment} algorithm applies: Alg.~3 groups terms by
relational scope and compares $\ind{P_1}+\cdots+\ind{P_n}$ against
$\ind{P'_1}+\cdots+\ind{P'_m}$ with $n\neq m$ allowed. This exact pattern is the paper's own
motivating example for that step.
\emph{Outside} $\mathcal F_T$ --- i.e.\ as soon as a \texttt{GroupBy}, \texttt{Distinct},
\texttt{Minus} or unknown operator appears --- \QED{} falls back to Alg.~6, which needs a
\textbf{bijection} between terms and returns \textsc{Unknown} when the counts differ. And
neither \QED{} algorithm relates sums with \emph{different relational scopes}; \liastar{}
does, and it applies uniformly rather than only inside a fragment.

\paragraph{(4) A sum of counts equals a count.}
\begin{lstlisting}[style=sql]
Q1: select a, count(b) from R group by a
Q2: select S.a, sum(S.cnt)
    from (select a, count(b) as cnt from R group by a) S group by S.a
\end{lstlisting}
\sqlsolver{} \textcolor{cGreen}{succeeds} --- it knows what \textsc{Sum} and \textsc{Count}
mean. \udp{} and \QED{} \textcolor{cRed}{cannot}, at any level of cleverness: with
aggregates uninterpreted, the required fact is simply not expressible. This is why
``aggregates'' is \QED{}'s largest Calcite failure bucket (96 cases).

\paragraph{(5) Redundant self-join on a unique key} (\cite{qed} \S6.3.2).
\begin{lstlisting}[style=sql]
Q1: select R.a from R join S S1 on R.a=S1.k join S S2 on R.a=S2.k
Q2: select R.a from R join S on R.a=S.k
\end{lstlisting}
\QED{} \textcolor{cGreen}{succeeds}: stabilisation finds $k_1=k_2=a$ through the SMT
congruence classes, then $\sq{S(a)}\times\sq{S(a)}=\sq{S(a)}$.
\spes{} \textcolor{cRed}{fails}: its self-join-on-key rule matches $R\bowtie(S\bowtie S)$, not
$(R\bowtie S)\bowtie S$, and it cannot infer $S_1.k=S_2.k$ transitively.
\textbf{Note this example has no top-level squash} --- the answer genuinely has multiplicity
2 --- so \QED{} discharges it by proving the summands equal \emph{pointwise}, not by turning
$\usum{}$ into $\exists$.

\paragraph{(6) Predicate pushdown through a left join} (\cite{qed} \S6.3.1).
\begin{lstlisting}[style=sql]
Q1: select * from (select * from R where R.a=42) X left join S on X.a=S.b
Q2: select * from (select * from R where R.a=42)
                  left join (select * from S where S.b=42) on TRUE
\end{lstlisting}
\QED{} \textcolor{cGreen}{succeeds}. \spes{} \textcolor{cRed}{fails}, and \QED{}'s paper
reports the reason as a soundness bug: it states that one of \spes{}'s normalisation rules
\emph{unsoundly} rewrites a left join whose condition is \texttt{TRUE} into a cross join.
(We report this as \cite{qed}'s claim; we did not reproduce it.) Anything needing NULL semantics
is out of reach for \udp{} entirely.

\paragraph{(7) \lstinline|ORDER BY| anything.} \sqlsolver{} only.

\subsection{The deep structural difference: squashed versus not}

Everything above has one root cause, and it is worth isolating because it predicts which
examples a given prover will handle.

When the top-level expression is \textbf{squashed} --- \lstinline|SELECT DISTINCT|,
\lstinline|GROUP BY|, \lstinline|EXISTS| --- multiplicities are irrelevant. \QED{} can then
push $\neq0$ inward, turning $\usum{}$ into $\exists$, $\times$ into $\wedge$ and $+$ into
$\vee$, and hand a pure Boolean formula to SMT. Cheap and strong.

When it is \textbf{not squashed} --- plain projection, a join, a \lstinline|UNION ALL| ---
the counts \emph{are} the answer, and the two systems fall back on comparing summands
pointwise, justified by
$(\forall \vec s.\ f(\vec s)=g(\vec s))\Rightarrow \usum{\vec s}f=\usum{\vec s}g$. That is
sound but demands \emph{alignable scopes}: \udp{}'s Alg.~2 and \QED{}'s Alg.~6 both require a
one-to-one match between terms. (\QED{}'s Alg.~3 relaxes this to ``same relational scope'',
which is strictly weaker but only available inside $\mathcal F_T$.)
\sqlsolver{} instead reasons about the sums' \emph{values}, which needs no alignment at all.

\keybox{Why the squash matters in practice}{%
Under a squash the comparison is Boolean, so a \lstinline|UNION ALL| split like
Example~\ref{ex:splitunion} reduces to
$\exists y.(y{\le}1000\wedge R{\neq}0)\vee\exists y.(y{>}1000\wedge z{<}500\wedge R{\neq}0)
\leftrightarrow \exists y.(\ldots)$ --- an ordinary SMT query that any SMT-backed tool settles.
Without the squash the same rewrite is an \emph{arithmetic} statement about sums. \QED{}
can still discharge it \emph{if} the terms share a relational scope and the query is in
$\mathcal F_T$ (Alg.~3 sums the bodies); otherwise it needs a term bijection.
\liastar{} needs neither condition.}

\subsection{On the published numbers}

\begin{table}[h]\centering\small
\begin{tabular}{@{}llrr@{}}
\toprule
reported by & benchmark described as & \spes{} & the paper's own tool\\
\midrule
\QED{}~\cite{qed}       & ``232 Calcite pairs'' & \textbf{95} & 147\\
\sqlsolver{}~\cite{sqlsolver} & ``232 Calcite pairs'' & \textbf{99} & 232\\
\bottomrule
\end{tabular}
\end{table}

\noindent The same baseline on nominally the same benchmark differs. The extraction
pipelines are different, so the suites are different, so the percentages are not
comparable. Additional cautions:

\begin{itemize}
\item \QED{}'s harder suites (444 non-trivial Calcite v1.32 pairs; 1287 CockroachDB pairs)
      were never run through \sqlsolver{}, and vice versa.
\item \QED{} does not evaluate against \sqlsolver{} at all --- the two are near-simultaneous.
\item \sqlsolver{}'s TPC-C/TPC-H pairs were generated by running Spark SQL's optimiser and
      then \emph{manually} converting logical plans back into SQL. That is a reasonable
      methodology but it is not an independent benchmark.
\item ``Proves 100\% of Calcite'' should be read as ``100\% of the pairs this group
      extracted and chose to model''.
\end{itemize}

\warnbox{If you are writing a paper in this area}{%
Do not cite a cross-paper number as a head-to-head comparison. The only defensible
comparison is one you run yourself, on one extraction, with all tools on the same input.
As of these two papers, nobody has done that.}

\subsection{What the combination would look like}

The two 2024 systems are complementary to an unusual degree, and the merge is not
speculative --- each piece has a clean interface.

\begin{center}
\begin{tikzpicture}[font=\footnotesize\sffamily,node distance=5mm,
  b/.style={draw=cGrey,fill=cGreyBg,rounded corners=3pt,inner sep=5pt,align=center,minimum height=9mm},
  q/.style={draw=cBlue,fill=cBlueBg,rounded corners=3pt,inner sep=5pt,align=center,minimum height=9mm},
  s/.style={draw=cAmber,fill=cAmberBg,rounded corners=3pt,inner sep=5pt,align=center,minimum height=9mm},
  ar/.style={-{Stealth[length=2.4mm]},cGrey,thick}]
\node[q] (m) {\QED{} modelling\\NULLs, PK/FK, QOp/HOp};
\node[q,right=of m] (st) {\QED{} stabilise\\SMT congruence};
\node[s,right=of st] (ag) {\sqlsolver{}\\concrete aggregates};
\node[s,right=of ag] (li) {\liastar{} backend\\no scope alignment needed};
\draw[ar](m)--(st);\draw[ar](st)--(ag);\draw[ar](ag)--(li);
\end{tikzpicture}
\end{center}

\QED{}'s term-bijection step is its ceiling; \liastar{} is exactly the tool that removes
that ceiling. \QED{}'s uninterpreted $\hop$ for aggregates is its largest failure bucket;
\sqlsolver{}'s aggregate tables fill it. Conversely \sqlsolver{}'s weakest areas --- NULLs,
constraints, unknown operators --- are precisely \QED{}'s contributions. No published system
does this.
